\section*{الفصل \ref{ch:g6:pure-substances-and-mixtures} --- المواد النقية والخلائط}

\begin{solution}{exo:g6:pure-substances-and-mixtures:1}
مواد نقية: الماء المقطر، وكتلة الحديد الخالص، وبلورة السكر. خلائط:
\omterm{def:g4:air-a-mixture-of-gases:air}{الهواء}، وماء البحر، والحليب.
\end{solution}

\begin{solution}{exo:g6:pure-substances-and-mixtures:2}
متجانسة: الماء المالح، وماء الصنبور. غير متجانسة: الغرانيت، والزيت فوق
الماء، والموسلي.
\end{solution}

\begin{solution}{exo:g6:pure-substances-and-mixtures:3}
صنف معيّن من المواد، له اسمه وخصائصه: مثل الماء، والملح، والإيثانول
(وأيضًا السكر، والأكسجين، والحديد).
\end{solution}

\begin{solution}{exo:g6:pure-substances-and-mixtures:4}
درجة غليان ثابتة توحي \omterm{def:g6:pure-substances-and-mixtures:pure-substance}{بمادة نقية}؛ و\qty{78}{\celsius} هي درجة غليان
الإيثانول.
\end{solution}

\begin{solution}{exo:g6:pure-substances-and-mixtures:5}
يحوي أنواعًا كيميائية كثيرة (ماء، وسكريات، وأحماض، وفيتامينات، ونكهات):
فهو \omterm{def:g6:pure-substances-and-mixtures:mixture}{خليط}. وكلمة «نقي» على العلبة تعني فقط أنه لم يُضف إليه
شيء.
\end{solution}

\begin{solution}{exo:g6:pure-substances-and-mixtures:6}
لا: \omterm{def:g6:pure-substances-and-mixtures:pure-substance}{المادة النقية} تغلي عند درجة حرارة ثابتة. إنه \omterm{def:g6:pure-substances-and-mixtures:mixture}{خليط}، على الأرجح ماء
ذابت فيه مادة ما، مثل الماء المالح.
\end{solution}

\begin{solution}{exo:g6:pure-substances-and-mixtures:7}
يستطيع المرشِّح أن يفصل الماء الموحل (تبقى الحبيبات على الورقة) والعصير
باللب (يبقى اللب). ولا يستطيع أن يفصل الماء المالح (فالملح ذائب)، ولا
الزيت عن الماء (فالسائلان كلاهما يمران عبره؛ ويُفصلان بصب الطبقة
العليا).
\end{solution}

\begin{solution}{exo:g6:pure-substances-and-mixtures:8}
$78.7 \div 10 = \qty{7.87}{g}$ لكل \qty{1}{cm^3}: إنه الحديد.
\end{solution}

\begin{solution}{exo:g6:pure-substances-and-mixtures:9}
$\frac{357}{478} \approx 0.747$: أي نحو \qty{75}{\percent}.
\end{solution}

\begin{solution}{exo:g6:pure-substances-and-mixtures:10}
$2 \times 35 = \qty{70}{g}$ من الأملاح. والنسبة المئوية: $\frac{35}{1000} =
\qty{3.5}{\percent}$.
\end{solution}

\begin{solution}{exo:g6:pure-substances-and-mixtures:11}
نقيس درجتي غليانهما (يجب أن تبقى كل منهما ثابتة في أثناء الغليان، وأن
تتساويا)، ونزن الحجم نفسه من كل منهما (كتلتان متساويتان). فثابتان
متساويان، كلاهما ثابت، يدلان على \omterm{def:g6:pure-substances-and-mixtures:pure-substance}{المادة النقية}
نفسها.
\end{solution}

\begin{solution}{exo:g6:pure-substances-and-mixtures:12}
$\frac{10}{90 + 10} = \qty{10}{\percent}$ من الإيثانول. إنه \omterm{def:g6:pure-substances-and-mixtures:mixture}{خليط}: فليس له
درجة غليان ثابتة ولا يرد في الجدول، الذي لا يذكر إلا المواد
النقية.
\end{solution}

\begin{solution}{pb:g6:pure-substances-and-mixtures:1}
\textbf{1.} لا. \omterm{def:g6:pure-substances-and-mixtures:pure-substance}{فالمادة النقية} \omterm{def:g6:pure-substances-and-mixtures:homogeneous}{والخليط المتجانس} قد يتطابقان في المظهر
تمامًا: لا بد من قياس شيء ما.

\textbf{2.} متجانس: فهو رائق وأجزاؤه لا تُرى.

\textbf{3.} ثابت: درجة غليان تبقى ثابتة في أثناء غليان السائل (أو كتلة
حجم معلوم، تُقارن بجدول).

\textbf{4.} A وB، اللذان تبقى درجتا غليانهما ثابتتين، مادتان نقيتان.
أما C \omterm{def:g6:pure-substances-and-mixtures:mixture}{فخليط}.

\textbf{5.} A يغلي عند \qty{100}{\celsius}: إنه الماء. B يغلي عند
\qty{78}{\celsius}: إنه الإيثانول.

\textbf{6.} A: $49.8 \div 50.0 = \qty{0.996}{g}$. B: $39.5 \div 50.0 =
\qty{0.790}{g}$. C: $51.2 \div 50.0 = \qty{1.024}{g}$.

\textbf{7.} نعم: نحو \qty{1.0}{g} للماء، و\qty{0.79}{g} للإيثانول.

\textbf{8.} كلما ذهب الماء بالغليان بقي الملح، فيحوي السائل الباقي ملحًا
أكثر فأكثر؛ فيغلي عندئذ عند درجة حرارة أعلى.

\textbf{9.} يحوي السائل C جسمًا صلبًا ذائبًا: فهو \omterm{def:g2:mixing-and-dissolving:dissolve}{محلول}، أي \omterm{def:g6:pure-substances-and-mixtures:mixture}{خليط}، كما بيّن
غليانه من قبل.

\textbf{10.} $100 \times 1.024 = \qty{102.4}{g}$؛ $\frac{3.5}{102.4}
\approx 0.034$: أي نحو \qty{3.4}{\percent} من الكتلة ملح.

\textbf{11.} $\frac{35}{1000} = \qty{3.5}{\percent}$: قريبة جدًا من
السؤال 10.

\textbf{12.} \qty{3.5}{g} في \qty{100}{mL}، إذن $3.5 \times 10 =
\qty{35}{g}$ من الملح في اللتر. فبملوحته وسلوكه عند الغليان، C هو ماء
البحر؛ وA هو الماء المقطر، وB هو الإيثانول.
\end{solution}
